\section{Retained arithmetic statements}\label{sec:retained49}
The following statements and their proofs later in the article are
retained from the preceding development, not consequences of the new
two-state alternative.

Our arithmetic conclusions go beyond uniform bad approximation. For $\alpha\in\R^r$ with $1,\alpha_1,\ldots,\alpha_r$ rationally independent, write
\begin{equation}\label{eq:type-intro}
 \omega^*(\alpha)=\limsup_{\norm m_1\to\infty}
       \frac{-\log\norm{m\cdot\alpha}_{\R/\Z}}{\log\norm m_1}.
\end{equation}
This is the ordinary dual exponent $\omega_{1,r}$ of the row vector $\alpha$ in \cite[Definition 1]{BL06}, with an equivalent fixed-dimensional norm; it is not the uniform exponent $\widehat\omega_{1,r}$ or the simultaneous exponent $\omega_{r,1}$. The alphabet consists of the identity and the rotations through $2\pi\alpha_i$. Throughout, $r$, the alphabet, $0<\rho\le1/10$ and
$0\le\varepsilon<\rho/(2\sqrt2)$ are fixed.

\begin{theorem}[Minimal Diophantine type]\label{thm:metric-main}
If $\omega^*(\alpha)=r$, then
\begin{equation}\label{eq:log-law}
 \lim_{N\to\infty}\frac{\log W_{N,\varepsilon}(\alpha)}{\log N}
                         =\frac{r}{2r+1}.
\end{equation}
For Lebesgue-almost every $\alpha$, and every fixed $\delta>0$, there are positive constants $c,C$, depending on the fixed $\alpha,\delta,r,\rho,\varepsilon$, such that, for $N\ge2$,
\begin{equation}\label{eq:ae-log-bounds}
 c\frac{N^{r/(2r+1)}}{(\log(N+2))^{2r(1+\delta)/(2r+1)}}
 \le W_{N,\varepsilon}(\alpha)
 \le C N^{r/(2r+1)}(\log(N+2))^{2r^2(1+\delta)/(2r+1)}.
\end{equation}
The upper inequalities are achieved by exact common-row machines. No uniformity near rational relations or in growing $r$ is asserted.
\end{theorem}

The exponent in \eqref{eq:log-law} need not imply a two-sided constant-factor law. The earlier dual badly approximable theorem supplies the stronger $\Theta$ conclusion on its stated class, and is recovered in Corollary~\ref{cor:ba}. Theorem~\ref{thm:metric-main} instead allows arbitrarily small Diophantine constants along subsequences. The quantitative arithmetic ingredients are classical; the memory conclusions concern their compatible realization and all-hidden-state converse.

There is also a genuine failure of a single exponent, even for one fixed irrational angle.
\begin{theorem}[Irrational horizons and Liouville fluctuations]\label{thm:liouville-main}
For every irrational $\alpha$ there are horizons $M_n\to\infty$ such that
$W_{M_n,\varepsilon}(\alpha)=\Theta(M_n^{1/3})$. If $\alpha$ is Liouville, then
\begin{equation}\label{eq:liouville-summary}
 \liminf_{N\to\infty}\frac{\log W_{N,\varepsilon}(\alpha)}{\log N}=0,
 \qquad
 \frac13\le\limsup_{N\to\infty}\frac{\log W_{N,\varepsilon}(\alpha)}{\log N}
 \le\frac12.
\end{equation}
Nevertheless $W_{N,\varepsilon}(\alpha)\to\infty$. All displayed upper bounds have exact realizations.
\end{theorem}

The upper limit in \eqref{eq:liouville-summary} is not determined here. The result establishes oscillation between subpower subsequences and a positive-power subsequence; it does not replace that distinction by a guessed exponent for all irrationals.


